\subsection{Incremental Entity Extraction}
\todo{maybe join with clustering}
\todo{incremental kbp}

As anticipated in \Cref{sec:incompleteness}, \acp{kr} and \acp{kb}\todo{ensure also kr are incomplete in sec:incompleteness} suffer from incompletess, and one of the use of \ac{ie} is to populate them (\acf{kbp})~\cite{entityorientedsearch2018balog}, for instance, with new movies extracted from newspapers.

The entities missing from a \ac{kr} are usually referred to as \emph{NIL entities}~\cite{entityorientedsearch2018balog}, where \emph{NIL} may stand for ``not in lexicon''~\cite{ILIEVSKI2020100565}. Another interpretation is that it simply refers to the English word ``nil'' that means ``nothing''~\cite{cambridge_nil}.

\begin{definition}[NIL Prediction]
    Given an entity mention $m$ and a reference \ac{kr}, \emph{NIL prediction} is the subtask of \acl{nel} that identifies whether $m$ referes to an entity that does not exist in the \ac{kr}, i.e., a \emph{NIL entity}~\cite{entity_linking_with_kb_2015}. 
\end{definition}
Through this work we will also speak of \emph{NIL mentions}, referring to mentions of NIL entities. 

This task has also been called with different names: ``unlinkable mention prediction''~\cite{entity_linking_with_kb_2015}, ``out-of-knowledge-Base mention discovery''~\cite{dong_reveal_the_unknown_2023}, ``out-of-knowledge-graph entity detection''~\cite{moller2024entity}, ``unknwon entity discovery''~\cite{kassner-etal-2022-edin}, and ``emerging entity discovery''~\cite{hoffart_discovering_emergin_entities_2014}.

\textcite{zhou-etal-2024-gendecider} distinguishes between NIL prediction and ``none of the candidates'' detection, arguing that \emph{NIL} means the referred entity is not in the \ac{kr}, while ``none of the candidates'' means the referred entity has not been retrieved without assuming it is not present in the \ac{kr}. However, in practise, NIL prediction is often performed by predicting whether the top-ranked entity, returned by the \ac{nel} ranking step, is correct or not~\cite{entity_linking_with_kb_2015}. It is impractical and prohibitive to compare a mention with every entity in a \ac{kr} for ensuring the mentions is NIL. Thus, in this work, we consider the preactical assumption that if a mention does not refer to the top-ranked candidate(s), returned by the \ac{nel} system, the mention is NIL. Under this assumption there is no difference between NIL prediction and none-of-the-candidates prediction.

\todo{definition of NIL clustering}
For addressing the incompleteness problem of \acp{kr}, once identified, NIL mentions that refer to the same NIL entity can be clustered to gather more contextual information, useful for representing the new entity. Indeed, \textcite{kassner-etal-2022-edin} achieved better results when constructing the entity representation from clusters of NIL mentions with respect to just using individual NIL mentions.

\begin{definition}[NIL Clustering]
    \emph{NIL clustering} is the task that aims at clustering NIL mentions that refer to the same NIL entity~\cite{ji2011tac}.
\end{definition}

NIL prediction has frequently been overlooked in \ac{nel}
systems, as many approaches exclude NIL mentions from their evaluation datasets and
eventually leave this task to future work~\cite{blink}. In fact, among the 38 approaches\todo{verify the published version}
compared by the survey \textcite{Sevgili2022NeuralEL} only 8 addressed NIL prediction. \todo{update for clustering too}

% dico che le TAC hanno dato importanza al NIL
Early work on NIL prediction dates back to 2006, with \textcite{bunescu-pasca-2006-using} using a threshold-based approach.
Later, the Text Analysis Conference (TAC)\footnote{\url{https://tac.nist.gov/}} fostered the research in NIL prediction by introducing NIL prediction in its \acf{kbp} track starting from 2009 \cite{mcnamee2009overview}. Starting from 2011, participants were also required to cluster together those NIL mentions that referred to the same entity \cite{ji2011tac}. Organized by NIST\footnote{\url{https://www.nist.gov/}}, TAC aimed to advance research in \ac{nlp} and related areas by offering large-scale test collections and standardized evaluation procedures. Since 2008, it has been held annually, inspiring a variety of studies on NIL prediction.

Also, from the 2015 edition, NIL prediction and NIL clustering have been included in the Named Entity rEcognition and Linking (NEEL) challenge~\cite{rizzo2017lessons}, a challenge focusing on microposts, i.e., posts from social media platforms such as x.com\footnote{\url{https://x.com/}}.


%% nil pred

\ac{nel} with NIL prediction can be framed as a classification problem with a reject option, which corresponds to NIL. According to \textcite{Sevgili2022NeuralEL}, four main strategies are typically used:
\begin{enumerate}
\item Assigning NIL whenever candidate generation produces no possible entities.
\item Applying a threshold to the linking score, below which the mention is treated as unlinkable.
\item Introducing NIL as an additional candidate during the ranking stage, effectively creating a new class in the classification problem.
\item Training a separate binary classifier on mention-entity pairs, often with extra features such as the top linking score or \ac{ner} labels, to decide whether a mention is linkable or not.
\end{enumerate}

Some studies \cite{Zhang2012TheNE, Rao2013EntityLF} argue that methods \ref{nilm-threshold} and \ref{nilm-add-class} are essentially equivalent, although the latter avoids manual threshold tuning.

%% nil approaches
Several approaches to NIL prediction have been proposed across the years.
Joint approaches to \ac{ner}, \ac{nel}, and NIL prediction have been proposed. \textcite{martins-etal-2019-joint} used an LSTM-based model to perform NER, NEL, and NIL prediction jointly. Earlier, \textcite{Fahrni2012HITSMA,Fahrni2013hits} jointly performed NEL, NIL prediction, and NIL clustering.

Threshold-based methods are also common. Systems in TAC 2017, such as those by \cite{BernierColborne2017CRIMsSF,Jiang2017SRCBED}, follow this approach, as do earlier works \cite{tan2015buptteam,Byrne2013UCDIA,huang2013msr,dalton2013umass}. These systems typically apply a score threshold to decide when a mention should be assigned NIL.

Another strategy is to add NIL as an additional class~\cite{Broscheit_2019,Rao2013EntityLF,Li2017AHM,klang2019overview,chen2013casia,barrena2013ubc,zhang2012nlprir,mcnamee2011cross}. For example, \textcite{Broscheit_2019} relied on BERT with a classification layer.

Binary classification approaches have also been widely studied~\cite{moreno_combining_word_entity_embeddings_2017,Yang2017TheTS,yang2013buptteam,Pink2013SYDNEYCA,zhou2013description,monahan2012lorify}. Some methods leveraged additional statistical or score-based features. \textcite{Yang2017TheTS}, for instance, incorporated the second-best candidate’s score and the standard deviation of the top-$k$ scores.

Finally, some systems rely on alternative methods. For example, \textcite{Paikens-EtAl:2017:TAC-KBP} used estimated coherence scores between the linked entities in the document.
\todo{add methods from moller}

% Advancements in NIL prediction continued also in parallel to the TAC conference~\cite{gottipati-jiang-2011-linking,hoffart_discovering_emergin_entities_2014,hoffart_awakens_2016,wu2016exploring}, even if \textcite{gottipati-jiang-2011-linking} experimentation was run on the TAC-KBP 2010 dataset~\cite{ji2010overview}. 

%%

More recently, since 2022, the interest in NIL prediction has motivated multiple studies to enable the popular \ac{nel} approach BLINK~\cite{blink} for NIL prediction~\cite{ijckg2022,kassner-etal-2022-edin,Heist2023nastylinker,zhu-etal-2023-learn}.

\textcite{ijckg2022}---that is one of the contributions of my Ph.D., and is better detailed in \Cref{ijckg2022}---proposed a benchmark dataset, an evaluation procedure, and a baseline pipeline for \emph{incremental \acl{nel}}. In this scenario, given $N$ batches of textual documents $b_0..b_N$, corresponding to $N$ time-steps, and an initial \ac{kr} $KR_0$, each time a batch of documents $b_k$ is processed, the identified NIL entities are added to the \ac{kr} $KR_{k+1}$, which will be used as the refernce \ac{kr} when processing the batch $b_{k+1}$.

Similarly, \textcite{kassner-etal-2022-edin} propose a dataset derived from Wikipedia and newspapers data, divided into two parts. The first part contains the documents exising before a time $t$, while the second the documents created after. The authors also proposed a pipeline-based system for incremental \acl{nel}.

Other datasets have been proposed for ...~\cite{nilk_dataset_2022,tempel_2022}.
\textcite{nilk_dataset_2022}. \textcite{tempel_2022} proposal is composed of ten yearly snapshots of Wikipedia from 2013 to 2022. NILK~\cite{nilk_dataset_2022} instead considers Wikidata as the reference \ac{kb}, using English Wikipedia for the  and is composed of two snapshots: Wikidata 2017 and 2022. 

... wikipedia wikidata

.... Wikipedia, ten snapshots from 2013 to 2022.

% main entity linking approach also include nil 
\cite{ayoola-etal-2022-refined}

% targeting NIL
\cite{nilk_dataset_2022,agarwal-etal-2022-entity,tempel_2022,,dong_reveal_the_unknown_2023,moller2024entity}

% llm based
\cite{zhou-etal-2024-gendecider,ding-etal-2024-chatel}

\todo{llm based}


\todo{neel challenge}

\todo{trattare insieme nil pred e clustering}

\todo{controlla chi fa incremental come noi in piu time steps. noi, edin}

\todo{moller paper reviewed a lot of approaches}
\todo{newer neural / bert based el systems do not target NIL pred, so contribution (better described later in section....)}
\todo{our contribution}





\todo{add recent studies}

In most cases, these techniques are incorporated into pipeline-based systems that treat NIL prediction and \ac{nel} as distinct steps~\cite{}. By contrast, other works have proposed joint models to reduce error propagation and leverage dependencies between the two tasks \cite{fahrni2013hits,martins-etal-2019-joint}.



\subsection{NIL Clustering}
\subsection{coreference resolution}
\todo{what did LLMs change?}
\subsection{Human-in-the-Loop}
